\documentclass[11pt,a4paper]{amsproc}

\usepackage[a4paper,margin=2.6cm]{geometry}
\usepackage[utf8]{inputenc}
\usepackage[T1]{fontenc}
\usepackage[spanish]{babel}   
\usepackage{amsmath,amssymb,amsthm}

\newtheorem*{theorem}{Teorema}

\newcommand{\C}{\mathbb{C}}

\begin{document}

\section*{Problema}

Dos aerolíneas, \textbf{Roja} y \textbf{Azul}, operan vuelos de ida entre $d$ ciudades;
cada una tiene su propio mapa de rutas (se permiten vuelos de una ciudad a sí misma).
Un \emph{itinerario} es una sucesión de $N$ tramos, uno por día, donde cada tramo se vuela
por una aerolínea determinada y por una ruta de esa aerolínea.

Para cada día $i$ el ministerio fija un par de precios $(\alpha_i,\beta_i)\neq(0,0)$ de
números complejos: $\alpha_i$ para un tramo Rojo de ese día y $\beta_i$ para un tramo Azul.
El costo de un itinerario es el \textbf{producto} de los $N$ precios que paga.

Probar que los precios se pueden fijar de antemano de modo que, para \textbf{todo} par
ordenado de ciudades $(u,v)$, el costo total de todos los itinerarios de $N$ días que van
de $u$ a $v$ sea exactamente cero. 

\section*{Solución}

La solución depende del siguiente resultado: 

\begin{theorem}
  Sean $A,B\in \C^{d\times d}$, $d\ge 1$. Entonces existen pares
$(\alpha_i,\beta_i)\neq(0,0)$, $1\le i\le N$, tales que
\[
  (\alpha_1A+\beta_1B)(\alpha_2A+\beta_2B)\cdots(\alpha_NA+\beta_NB)=0 .
\]
Más aún, se puede tomar $N\le 2^{d}-1$.
\end{theorem}

Sean $A,B\in\C^{d\times d}$ las matrices de adyacencia de Roja y Azul, es decir
$A_{xy}=1$ si Roja vuela $x\to y$ y $A_{xy}=0$ si no, e igual para $B$ con Azul.
El ``mapa tarifado'' del día $i$ es $C_i=\alpha_iA+\beta_iB$, cuya entrada $(x,y)$ es
justamente el precio del tramo $x\to y$ ese día (sumado sobre las aerolíneas que lo vuelan).
Multiplicando,
\[
  (C_1C_2\cdots C_N)_{uv}
  \;=\!\!\sum_{u=x_0\to x_1\to\cdots\to x_N=v}\ \prod_{i=1}^N (C_i)_{x_{i-1}x_i}
  \;=\!\!\sum_{\text{itinerarios } u\to v}\!\!\text{costo},
\]
porque desarrollar el producto de las entradas equivale a elegir, para cada día, una ruta
\emph{y} la aerolínea que la vuela. Por lo tanto la conclusión pedida es exactamente
$C_1C_2\cdots C_N=0$, que es el Teorema aplicado a $A$ y $B$. \qedhere

\section*{Observación}

Ll problema es imposible con precios reales y resoluble con 
precios complejos. El teorema usa fuertemente que
uno está en los complejos. 

Sea $d=2$, con Roja $=$ un bucle en cada ciudad más la ruta $1\to 2$, y Azul $=$ un bucle
en cada ciudad más la ruta $2\to 1$:
\[
  A=\begin{pmatrix}1&1\\0&1\end{pmatrix},\qquad
  B=\begin{pmatrix}1&0\\1&1\end{pmatrix},\qquad
  \alpha A+\beta B=\begin{pmatrix}\alpha+\beta&\alpha\\ \beta&\alpha+\beta\end{pmatrix},
\]
\[
  \det(\alpha A+\beta B)=(\alpha+\beta)^2-\alpha\beta=\alpha^2+\alpha\beta+\beta^2
  =\Big(\alpha+\tfrac{\beta}{2}\Big)^{\!2}+\tfrac34\,\beta^2 .
\]
Con precios \emph{reales} esto es $>0$ para todo $(\alpha,\beta)\neq(0,0)$: la matriz de cada
día es inversible, luego el producto también lo es, y los itinerarios nunca se cancelan,
cualquiera sea $N$. En cambio, con $\omega=e^{2\pi i/3}$ (o sea $1+\omega+\omega^2=0$) la
forma se anula y bastan tres días:
\[
  \bigl(\omega A+B\bigr)\bigl((3+\omega)A+(\omega-2)B\bigr)\bigl(\omega A+B\bigr)=0 .
\]
\end{document}
